\begin{proof}[Proof of \cref{obj:lem:fixed-released-law-baseline}]
\noindent\emph{1. Construct compatible endpoint laws.}
Compatibility implies that \(H\) and \(P_{\mathrm{rel}}\) are probability laws and, for every arm \(a\) and label \(r\),
\[
P_{\mathrm{rel}}(A=a,R=r)=q_{ar}.
\]
By \cref{obj:ass:overlap}, \(\varepsilon\le p_a(e)\le 1-\varepsilon\) on \(\mathcal E\). Thus
\[
\varepsilon H(C_r)\le q_{ar}\le(1-\varepsilon)H(C_r),
\]
so a zero-mass arm cell is also \(H\)-null. Define, on \(\mathcal E\),
\[
w_a(e)=\frac{1}{p_a(e)}.
\]
For each \(q_{ar}>0\), let \(U\) be uniform on \((0,1)\) and define two laws on outcome–weight pairs:
\[
\pi_{ar}^{-}
=\mathcal L\bigl(Q_{F_{ar}}(U),Q_{W_{ar}}(1-U)\bigr),
\qquad
\pi_{ar}^{+}
=\mathcal L\bigl(Q_{F_{ar}}(U),Q_{W_{ar}}(U)\bigr).
\]
Both have marginals \(F_{ar}\) and \(\mathcal L(W_{ar})\). They can be lifted to laws \(\gamma_{ar}^{s}\) on score–outcome pairs, for \(s\in\{-,+\}\), satisfying
\[
\mathcal L_{\gamma_{ar}^{s}}(E)=G_{ar},
\qquad
\mathcal L_{\gamma_{ar}^{s}}\bigl(Y,w_a(E)\bigr)=\pi_{ar}^{s}.
\]
Indeed, \(\mathcal L(W_{ar})\) is the common weight marginal of \(\pi_{ar}^{s}\) and \(\mathcal L(E_{ar},w_a(E_{ar}))\). Conditional on that weight, couple their outcome and score coordinates. Conditional laws exist on these standard Borel spaces.

For either sign, form the assigned and potential-outcome laws
\[
N_a^{s}=\sum_{r:q_{ar}>0}q_{ar}\gamma_{ar}^{s},
\qquad
M_a^{s}(de,dy)=w_a(e)\,N_a^{s}(de,dy).
\]
The score marginal of \(N_a^{s}\) is \(p_a(e)H(de)\): on \(C_r\), its score marginal is \(q_{ar}G_{ar}(de)=\mathbf1_{C_r}(e)p_a(e)H(de)\), and zero-mass cells are \(H\)-null. Consequently,
\[
M_a^{s}(de,[0,1])=H(de),
\qquad
p_a(e)M_a^{s}(de,dy)=N_a^{s}(de,dy).
\]
For signs \(s_0,s_1\in\{-,+\}\), glue \(M_0^{s_0}\) and \(M_1^{s_1}\) along their common score marginal \(H\), obtaining a law of \((E,Y_0,Y_1)\) with those two score–outcome marginals. Conditional on the resulting variables, draw \(A=1\) with probability \(E\), and set \(R=g(E)\) and \(Y=Y_A\). Denote the resulting full law by \(P^{s_0s_1}\).

This construction has score law \(H\), score-based assignment, consistency, and deterministic release. For \(q_{ar}>0\) and measurable \(B\subseteq[0,1]\), it also gives
\[
P^{s_0s_1}(R=r,A=a,Y\in B)
=\int_{C_r\times B}p_a(e)\,M_a^{s_a}(de,dy)
=q_{ar}F_{ar}(B)
=P_{\mathrm{rel}}(R=r,A=a,Y\in B).
\]
Both sides vanish when \(q_{ar}=0\). Hence every \(P^{s_0s_1}\) belongs to the class in \cref{obj:def:compatible-law-class}. % lean: CausalSmith.PartialID.UnlinkedPropensityAte.endpointFullLaw_compatibleCausalLaw

\noindent\emph{2. Compute their means and treatment effects.}
The two quantile couplings give, for every positive-mass cell,
\[
\int yw\,\pi_{ar}^{-}(dy,dw)
=\int_0^1 Q_{F_{ar}}(u)Q_{W_{ar}}(1-u)\,du,
\qquad
\int yw\,\pi_{ar}^{+}(dy,dw)
=\int_0^1 Q_{F_{ar}}(u)Q_{W_{ar}}(u)\,du.
\]
Since \(M_a^{s}=w_aN_a^{s}\), summing these identities over cells yields
\[
\mathbb E_{P^{s_0s_1}}[Y_a]
=\sum_{r:q_{ar}>0}q_{ar}\int yw\,\pi_{ar}^{s_a}(dy,dw)
=
\begin{cases}
\underline\mu_a,&s_a=-,\\
\overline\mu_a,&s_a=+.
\end{cases}
\]
Set
\[
P_{\mathrm{lo}}=P^{--},\qquad P_{\mathrm{hi}}=P^{++},
\qquad P^-=P^{+-},\qquad P^+=P^{-+},
\]
where the first sign selects the control-arm coupling. Bounded outcomes permit subtraction of their expectations, so
\begin{equation}\label{eq:fixed-law-ate-endpoints}
\mathbb E_{P^-}[Y_1-Y_0]=\underline\mu_1-\overline\mu_0,
\qquad
\mathbb E_{P^+}[Y_1-Y_0]=\overline\mu_1-\underline\mu_0.
\end{equation}
% lean: CausalSmith.PartialID.UnlinkedPropensityAte.endpointFullLaw_armMean

\noindent\emph{3. Bound every compatible arm mean.}
Fix \(P\in\mathfrak M(H,g,P_{\mathrm{rel}})\). For \(q_{ar}>0\), define its conditional outcome–weight coupling by
\[
\pi_{ar,P}
=\mathcal L_P\bigl(Y_a,w_a(E)\mid A=a,R=r\bigr).
\]
Consistency and the released-law condition give its outcome marginal \(F_{ar}\); score-based assignment and score marginality give its weight marginal \(\mathcal L(W_{ar})\). The same assignment identity, now weighted by \(w_a(E)\), gives
\[
\mathbb E_P[Y_a]
=\sum_{r:q_{ar}>0}q_{ar}\int yw\,\pi_{ar,P}(dy,dw).
\]
Cells with \(q_{ar}=0\) contribute zero.

For any coupling \(\pi\) of \(F_{ar}\) and \(\mathcal L(W_{ar})\), the layer-cake identity writes \(\int yw\,d\pi\) as the integral of joint upper-tail probabilities. For thresholds \(\tau_Y,\tau_W\ge0\), those probabilities lie between the intersection bounds
\[
\max\!\left\{P(Y>\tau_Y)+P(W>\tau_W)-1,0\right\}
\le P(Y>\tau_Y,W>\tau_W)
\le \min\!\left\{P(Y>\tau_Y),P(W>\tau_W)\right\}.
\]
The couplings \(\pi_{ar}^{-}\) and \(\pi_{ar}^{+}\) attain the respective bounds: their quantile upper-tail events are oppositely ordered and nested, respectively. Integrating the tail bounds therefore gives
\[
\int_0^1 Q_{F_{ar}}(u)Q_{W_{ar}}(1-u)\,du
\le \int yw\,\pi_{ar,P}(dy,dw)
\le \int_0^1 Q_{F_{ar}}(u)Q_{W_{ar}}(u)\,du.
\]
All integrals are finite because \(0\le Y_a\le1\) and \(w_a(E)\le\varepsilon^{-1}\). Multiplication by \(q_{ar}\) and summation show that
\begin{equation}\label{eq:fixed-law-arm-bound}
\underline\mu_a\le\mathbb E_P[Y_a]\le\overline\mu_a.
\end{equation}
% lean: CausalSmith.PartialID.UnlinkedPropensityAte.compatible_armMean_mem_endpoints

\noindent\emph{4. Fill each arm interval.}
For any \(x\in[\underline\mu_a,\overline\mu_a]\), choose \(\rho_a\in[0,1]\) with
\[
x=(1-\rho_a)\underline\mu_a+\rho_a\overline\mu_a,
\]
including \(\rho_a=0\) when the endpoints coincide, and define
\[
P_{a,x}=(1-\rho_a)P_{\mathrm{lo}}+\rho_aP_{\mathrm{hi}}.
\]
Probability normalization, the score and released marginals, and the assignment identity are preserved by this mixture; consistency and deterministic release hold under both component laws and hence under the mixture. Thus \(P_{a,x}\) is compatible. Linearity of integration gives
\[
\mathbb E_{P_{a,x}}[Y_a]
=(1-\rho_a)\mathbb E_{P_{\mathrm{lo}}}[Y_a]
+\rho_a\mathbb E_{P_{\mathrm{hi}}}[Y_a]=x.
\]
Together with \cref{eq:fixed-law-arm-bound}, this proves the claimed attainable interval for each arm. % lean: CausalSmith.PartialID.UnlinkedPropensityAte.compatibleCausalLaw_mixture

\noindent\emph{5. Fill the treatment-effect interval.}
For every compatible \(P\), the two arm bounds and
\(\mathbb E_P[Y_1-Y_0]=\mathbb E_P[Y_1]-\mathbb E_P[Y_0]\) give
\[
\underline\mu_1-\overline\mu_0
\le\mathbb E_P[Y_1-Y_0]
\le\overline\mu_1-\underline\mu_0.
\]
Conversely, for any \(x\) in this interval, choose \(\rho_{\mathrm{ATE}}\in[0,1]\) such that
\[
x=(1-\rho_{\mathrm{ATE}})(\underline\mu_1-\overline\mu_0)
+\rho_{\mathrm{ATE}}(\overline\mu_1-\underline\mu_0),
\]
and set
\[
P_x=(1-\rho_{\mathrm{ATE}})P^-+\rho_{\mathrm{ATE}}P^+.
\]
The same mixture argument makes \(P_x\) compatible, and linearity gives
\(\mathbb E_{P_x}[Y_1-Y_0]=x\). The interval is \(I(P_{\mathrm{rel}};H,g)\) by \cref{obj:def:sharp-ate-set}; the compatible laws \(P^-\) and \(P^+\) attain its two endpoints as computed in \cref{eq:fixed-law-ate-endpoints}. % lean: CausalSmith.PartialID.UnlinkedPropensityAte.fixedReleasedLaw_sharp
\end{proof}
